% =============================================================================
% Section 5: Robustness (Short)
% =============================================================================
\section{Robustness}\label{sec:robustness}

The duplex finding must survive interrogation.  This section reports
three robustness checks drawn from Paper~2~\cite{isi_paper2}.

\subsection{Defence omission stress test}

The LOO analysis in \cref{sec:core-finding} removes each axis
individually.  The defence-omission variant yields
$\rho = 0.833$---the lowest correlation with the baseline among all
six single-axis omissions.  The mean absolute rank displacement is
3.33~positions; the maximum is 10~positions.

For comparison, the logistics-omission variant yields
$\rho = 0.926$; critical inputs, $\rho = 0.941$; technology,
$\rho = 0.968$; energy, $\rho = 0.992$; financial, $\rho = 0.993$.
The ordering is monotonic in variance share: axes that contribute
more variance produce larger rank disruptions when removed.

\subsection{Dirichlet weight perturbation}

The equal-weight assumption ($w_j = 1/6$ for all~$j$) is a design
choice.  To test sensitivity, we draw 10\,000 weight vectors from
a symmetric Dirichlet distribution with $\alpha_j = 10$, recompute
the composite as $C_i = \sum_j w_j A_{ij}$, and record the
resulting ranking.

The mean rank-order Spearman correlation across all 10\,000 draws
is $\rho = 0.980$.  The ranking is structurally stable under
moderate weight perturbation.  The highest rank volatility is
observed for the Netherlands ($\sigma = 3.08$), a country whose
defence score (0.960) is exceptionally high relative to its
composite rank (15th).

\Cref{tab:dirichlet-top} reports the eight most volatile countries.

% Table T05: Dirichlet weight perturbation -- highest volatility countries
\begin{table}[htbp]
\centering
\caption{Dirichlet weight perturbation ($\alpha_j=10$, $N_{\text{MC}}=10{,}000$, seed~42): eight highest-volatility countries by rank standard deviation.  Aggregate mean Spearman~$\rho = 0.980$.}
\label{tab:dirichlet-top}
\begin{tabular}{@{}r l rrrr@{}}
\toprule
\textbf{Base Rk} & \textbf{Country} & \textbf{Mean Rk} & $\boldsymbol{\sigma}$ & \textbf{P5} & \textbf{P95} \\
\midrule
  15 & Netherlands & 15.35 & 3.08 & 10.0 & 20.0 \\
  6 & Finland & 6.44 & 2.38 & 4.0 & 12.0 \\
  13 & Estonia & 13.50 & 2.21 & 9.0 & 17.0 \\
  9 & Italy & 9.46 & 2.17 & 6.0 & 13.0 \\
  14 & Bulgaria & 14.12 & 2.02 & 11.0 & 18.0 \\
  11 & Greece & 10.75 & 1.93 & 8.0 & 14.0 \\
  5 & Croatia & 5.84 & 1.75 & 3.0 & 9.0 \\
  7 & Sweden & 6.31 & 1.73 & 3.0 & 9.0 \\
\bottomrule
\end{tabular}
\end{table}

The Netherlands result deserves comment.  Its defence score is the
third-highest in the sample (behind Malta and Slovakia's zero), yet
its scores on all other axes are low to moderate.  When weight draws
increase the defence weight, the Netherlands rises; when they
decrease it, it falls.  The country's rank volatility is a direct
consequence of its extreme axis profile---and a microcosm of the
duplex: the defence axis determines where the Netherlands sits.

\subsection{Correlation structure}

The variance decomposition in \cref{sec:proof} depends on the
covariance matrix of the six axes.  If axes were strongly
correlated, the cross-covariance terms would be large and the
decomposition less informative.  In practice, the correlation
structure is sparse.

Of 15~pairwise Pearson correlations among the six axes, only two
are statistically significant at $\alpha = 0.05$ after Bonferroni
correction (15~tests, critical $|t| \approx 2.79$, $\text{df} = 25$):

\begin{itemize}[leftmargin=2em]
  \item Energy--Technology: $t = 3.12$.
  \item Critical Inputs--Logistics: $t = 3.28$.
\end{itemize}

\noindent
The financial axis is orthogonal to all others: its maximum
$|t|$-statistic against any axis is 1.11.  This confirms that
financial concentration measures a structurally independent
dimension of external dependence---one that happens to exhibit
minimal cross-country dispersion.

\begin{sloppypar}
The sparse correlation structure reinforces the variance
decomposition result: the axes are largely independent, so the
marginal contributions reflect genuine structural content rather
than statistical redundancy.
\end{sloppypar}
